% !TeX root = ../papers/nc-p3.tex
\section{正の固定曲線の端点・相方・八次閉点の全次数認証}
\label{sec:fixed-residual-regularity}
\begin{proposition}\label{prop:fixed-residual-regularity}
倍率 $(z,\lambda)=(3,2)$ の正の固定曲線上で、次の全22幾何点は反対階数16の
標準Koszul左右AS正則窓を与える。
\begin{enumerate}
\item $a=0$ と $b=0$ の二つの許容点。
\item $H_\epsilon=0$ の四つの第一座標の上で、$\det A+\det B=0$ を満たす点の相方。
\item 命題\ref{prop:fixed-quintic-slice}の小行列式の消去に残る二つの八次閉点の全共役。
\end{enumerate}
各始点の道空間の直和は四つの中心四次式の多項式環上で階数256の自由加群となる。
最初の二項の六点は、$(3,2)$ を含む倍率平面の共通の非空開集合上で同じ性質を保つ。
\end{proposition}
\begin{proof}
定理\ref{thm:uniform-fixed-quartic}の端点と相方の有理式を $r^2=3$ で読む。
相方の第一座標は二つの相異なる二次式の根であり、二つの端点から離れている。
193で $r$ は $14,179$ に単純分解する。この二値で相方の式を、端点では一方の値で式を還元する。
全ての分母と $xy\det A\det B(\det A+\det B)$ は剰余体で非零である。

八次閉点は同命題の検算が返す $(P(a),u-U(a))$ を重複なく取ると二つであり、
第一座標の既約八次式 $P$ も相異なる。各 $P$ の一つの単根を1009で取り、
$(a,U(a))$ の還元は $(60,14)$ と $(33,982)$ となる。
$P$ と $U$ の係数の分母はこの素数で可逆である。
Henselの補題で単根は持ち上がり、$\mathbb Q[a]/(P)$ の一つの付値で窓を還元できる。
証明すべき性質は体拡大と共役で保たれるので、各八次式につき一つの還元で足りる。

\path{computations/verify_fixed_residual_regularity.py} はこの八つの剰余窓で四つの関係空間を
作り、隣接階数6、反対階数16、四始点の長さ4,5の次元35,56と
四次中心の文字次元 $(3,2,0,0)$ を検算する。最初の二文字から各二式を選んだ商の次元は
全32組の窓と始点で、次数0から13まで順に
\[
 1,4,10,20,31,40,44,40,31,20,10,4,1,0
\]
となる。定理\ref{thm:fixed-line-elliptic}の持上げの議論を適用する。
すなわち経路と中心の次元は族の一般点から来る下界に等しく、
四中心は付値環上に持ち上がる。剰余商の次元上界と道空間の全次数の下界から
自由基底と左右Koszul完全性が従う。端点と相方の表示では分母の非消滅、経路と中心の階数、
商の有限性の証明に用いた小行列式がそれぞれ開近傍で保たれる。
$r^2=z$ の有限被覆で六点の近傍の補集合の像を除けば、所要の倍率の開集合を得る。
\end{proof}
